\documentclass[a4paper, serif, 10pt]{moderncv}

%% moderncv
% color and style
\moderncvcolor{black}
\moderncvstyle{banking}

% renew moderncv command to adapt for CJK colon
\renewcommand*{\cvitem}[3][.25em]{%
  \ifstrempty{#2}{}{\hintstyle{#2}：}{#3}%
  \par\addvspace{#1}}

\renewcommand*{\cvitemwithcomment}[4][.25em]{%
  \savebox{\cvitemwithcommentmainbox}{\ifstrempty{#2}{}{\hintstyle{#2}：}#3}%
  \setlength{\cvitemwithcommentmainlength}{\widthof{\usebox{\cvitemwithcommentmainbox}}}%
  \setlength{\cvitemwithcommentcommentlength}{\maincolumnwidth-\separatorcolumnwidth-\cvitemwithcommentmainlength}%
  \begin{minipage}[t]{\cvitemwithcommentmainlength}\usebox{\cvitemwithcommentmainbox}\end{minipage}%
  \hfill% fill of \separatorcolumnwidth
  \begin{minipage}[t]{\cvitemwithcommentcommentlength}\raggedleft\small\itshape#4\end{minipage}%
  \par\addvspace{#1}}

%% page layout/margins
\usepackage[top=2.5cm, bottom=2.5cm, left=1.5cm, right=1.5cm]{geometry}

\nopagenumbers{}

%% Babel config for English language
\usepackage[english]{babel}

%% fontspec
\usepackage{fontspec}

\IfFontExistsTF{Linux Libertine}{
  \setmainfont[Ligatures={TeX, Common}, Numbers=Lining, ItalicFont=Linux Libertine]{Linux Libertine}
}{}
\IfFontExistsTF{Linux Libertine O}{
  \setmainfont[Ligatures={TeX, Common}, Numbers=Lining, ItalicFont=Linux Libertine O]{Linux Libertine O}
}{}

%% CTeX
% CJK support, used to show CJK characters in the resume
%
% - heading=false: disable ctex heading
% - punct=kaiming: use kaiming punctuations styles for CJK
% - scheme=plain: use plain scheme, do not override \normalsize font size
% - space=auto: space settings for CJK characters
%
% CJK fonts are resolved per locale language, e.g. Noto Serif/Sans CJK JP
% for Japanese resumes, with fallback to the next region if not installed
%
% ref:
% - http://ctan.mirrorcatalogs.com/language/chinese/ctex/ctex.pdf
\usepackage[UTF8, heading=false, punct=kaiming, scheme=plain, space=auto]{ctex}

\IfFontExistsTF{Noto Serif CJK JP}{
  \setCJKmainfont{Noto Serif CJK JP}
}{}
\IfFontExistsTF{Noto Sans CJK JP}{
  \setCJKsansfont{Noto Sans CJK JP}
}{}

%% line spacing
\usepackage{setspace}
\setstretch{1.125}

%% URL styling - use normal text instead of monospace
\urlstyle{same}

%% Patch \httplink and \httpslink to handle full URLs
% The original moderncv macros always prepend http:// or https:// to URLs.
% This causes issues when the URL already has a protocol (e.g., https://example.com)
% resulting in malformed URLs like https://https://example.com.
% This patch checks if the URL already contains "://" and uses it directly if so.
% Note: We use \str_set:Nx (x-expansion) to expand macros like \@homepage first.
\ExplSyntaxOn
\str_new:N \g_yamlresume_url_str

\RenewDocumentCommand{\httpslink}{O{}m}{
  \str_set:Nx \g_yamlresume_url_str {#2}
  \str_if_in:NnTF \g_yamlresume_url_str {://}
    { \url{#2} }
    { \url{https://#2} }
}

\RenewDocumentCommand{\httplink}{O{}m}{
  \str_set:Nx \g_yamlresume_url_str {#2}
  \str_if_in:NnTF \g_yamlresume_url_str {://}
    { \url{#2} }
    { \url{http://#2} }
}
\ExplSyntaxOff

\name{佐藤 健太}{}
\title{データエンジニア / データ基盤アーキテクト}
\phone[mobile]{+81 90-1234-5678}
\email{kenta.sato@example.jp}
\homepage{https://kenta-sato.dev}

\address{日本、東京都、東京、東京都渋谷区神南1-2-3、150-0041}{}{}

\extrainfo{{\small \faGithub}\ \href{https://github.com/kenta-data}{@kenta-data} {} {} {} • {} {} {} 
{\small \faLinkedin}\ \href{https://www.linkedin.com/in/kenta-sato-data}{@kenta-sato-data}}

\begin{document}

\maketitle

\section{基本情報}

\cvline{}{データエンジニアとして8年以上の経験を持ち、大規模データ基盤の設計・構築・運用を一貫して担当してきました。ビジネス要件を深く理解し、信頼性と拡張性の高いデータパイプラインを提供することで、データドリブンな意思決定を支援します。
%
\begin{itemize}
\item クラウド（AWS/GCP）を活用したデータレイク・データウェアハウスの構築と運用
\item Apache Spark、Airflow、dbt を用いたバッチ/ストリーム処理の設計・実装
\item データガバナンス、品質管理、コスト最適化の推進
\item クロスファンクショナルなチームでのリードとメンタリング
\end{itemize}}
\section{学歴}

\cventry{2011年4月～2015年3月}
        {学士、情報理工学、成績：3.7}
        {東京大学}
        {\url{https://www.u-tokyo.ac.jp}}
        {}
        {\begin{itemize}
\item 情報理工学を専攻し、データベース、分散システム、機械学習の基礎を習得
\item 卒業研究では大規模データ処理の効率化をテーマに研究を実施
\item チーム開発やプレゼンテーションを通じて協働スキルを向上
\end{itemize}
\textbf{コース}：データ構造とアルゴリズム、データベースシステム、分散システム、機械学習、統計学、情報理論}
\section{職歴}

\cventry{2021年4月～現在}
        {シニアデータエンジニア}
        {株式会社データフォージ}
        {\url{https://www.dataforge.example.jp}}
        {}
        {\begin{itemize}
\item データ基盤のアーキテクチャ設計をリードし、Snowflake と dbt を中心とした分析基盤を構築
\item Airflow による数百本のデータパイプラインのオーケストレーションを設計・運用
\item データ品質監視の仕組みを導入し、重要指標の信頼性を向上
\item ジュニアエンジニアのメンタリングとコードレビューを実施
\end{itemize}
\textbf{キーワード}：データ基盤、Snowflake、dbt、Airflow、データ品質}

\cventry{2018年4月～2021年3月}
        {データエンジニア}
        {株式会社ネクストリテイル}
        {\url{https://www.nextretail.example.jp}}
        {}
        {\begin{itemize}
\item AWS 上にデータレイクを構築し、POS や EC サイトのデータを統合
\item Apache Spark を用いた大規模バッチ処理の設計・実装を担当
\item データカタログを整備し、社内のデータ探索と利用を促進
\item コスト最適化のため、S3 ライフサイクルポリシーと EMR の実行計画を改善
\end{itemize}
\textbf{キーワード}：AWS、Apache Spark、データレイク、データカタログ、コスト最適化}

\cventry{2015年4月～2018年3月}
        {データ分析エンジニア}
        {株式会社インフォメトリクス}
        {\url{https://www.infometrics.example.jp}}
        {}
        {\begin{itemize}
\item 受託分析プロジェクトにて、データ抽出・加工・可視化を担当
\item SQL と Python を用いた分析基盤の構築とレポーティングの自動化
\item 顧客向けの KPI ダッシュボードを Tableau で設計・開発
\item 要件定義から運用までを一貫して経験
\end{itemize}
\textbf{キーワード}：Python、SQL、Tableau、自動化、KPI}
\section{言語}

\cvline{日本語}{ネイティブ・母国語レベル \hfill \textbf{キーワード}：日本語ネイティブ}
\cvline{英語}{完全な実務レベル \hfill \textbf{キーワード}：TOEIC 900、技術文書読解}
\cvline{中国語}{限定的な実務レベル \hfill \textbf{キーワード}：HSK 4級}
\section{スキル}

\cvline{データエンジニアリング}{エキスパート \hfill \textbf{キーワード}：Python、SQL、Apache Spark、Airflow、dbt、Kafka}
\cvline{クラウド \& インフラ}{上級 \hfill \textbf{キーワード}：AWS、GCP、Terraform、Docker、Kubernetes}
\cvline{データウェアハウス}{上級 \hfill \textbf{キーワード}：Snowflake、BigQuery、Redshift、PostgreSQL}
\cvline{分析 \& 可視化}{中級 \hfill \textbf{キーワード}：Tableau、Looker、Pandas、Jupyter}
\section{受賞・表彰}

\cventry{2023年10月}
        {株式会社データフォージ}
        {技術表彰 最優秀賞}
        {}
        {}
        {データ基盤の刷新プロジェクトにおいて、設計・実装・運用面で顕著な成果を上げたとして表彰。}
\section{資格・免状}

\cventry{2022年3月}
        {Amazon Web Services}
        {AWS Certified Data Analytics - Specialty}
        {\url{https://aws.amazon.com/certification/certified-data-analytics-specialty/}}
        {}
        {}

\cventry{2020年6月}
        {Google Cloud}
        {Google Cloud Professional Data Engineer}
        {\url{https://cloud.google.com/certification/data-engineer}}
        {}
        {}
\section{出版・著作}

\cventry{2022年11月}
        {データ基盤の内製化におけるAirflow活用事例}
        {技術評論社}
        {\url{https://gihyo.jp/article/data-platform-airflow}}
        {}
        {\begin{itemize}
\item データ基盤の内製化を進める際のAirflow活用方法を解説
\item パイプライン設計、運用監視、チーム体制づくりに焦点
\item 実案件での課題と解決策を紹介
\end{itemize}}
\section{プロジェクト}

\cventry{2022年1月～2022年12月}
        {顧客データを統合し、マーケティングとサービス改善を支援するデータ基盤}
        {顧客360度データ基盤の構築}
        {\url{https://www.dataforge.example.jp/case/customer360}}
        {}
        {\begin{itemize}
\item 複数システムの顧客データをSnowflakeに統合
\item dbtによるデータモデル整備と品質テストの自動化
\item BIツール向けのセマンティックレイヤーを設計
\item マーケティング部門と連携し、施策立案のリードタイムを短縮
\end{itemize}
\textbf{キーワード}：顧客データ、Snowflake、dbt、データモデリング}
\section{興味・関心}

\cvline{オープンソース}{Airflow、dbt、コミュニティ活動}
\cvline{読書}{技術書、データアーキテクチャ}
\section{ボランティア活動}

\cventry{2020年1月～現在}
        {メンター}
        {データエンジニア勉強会コミュニティ}
        {\url{https://data-engineer-community.example.jp}}
        {}
        {\begin{itemize}
\item データエンジニアを目指す参加者向けに勉強会やハンズオンを実施
\item キャリア相談や技術的な質問への回答を通じて成長を支援
\item コミュニティ運営にも携わり、継続的な学びの場を提供
\end{itemize}}
    
\end{document}